% CPSY 1291 -- Recitation handout 11
% Compile:  latexmk -lualatex handout-11-attention-transformers-generative.tex
\documentclass[11pt]{article}
\usepackage{cpsy1291-handout}

\begin{document}

\handouthead{11}{Attention, transformers, and generative models}
  {Three lectures of mechanics, in one session}
  {Week 11 \ $\cdot$\ hold this after Lecture 18; Assignment 5 is due Tue 11/24}

\section*{What this session is for}
\addcontentsline{toc}{section}{What this session is for}

This is the longest handout of the term and the last one about content, because
the calendar leaves no alternative: Assignment 5 covers Lectures 15--18, it is
due on Tuesday 11/24, Thanksgiving recess begins the day after, and the final
exam is the first day back. So Lectures 16, 17 and 18 have to share one session.

Read it in three passes rather than one. Section~1 (attention) is the part
everything else depends on, and it is worth doing slowly with a pencil. Section~2
(generation) is short. Section~3 (VAEs and diffusion) is two objectives, each of
which is five lines of code once you have seen it written down.

\begin{keybox}
\ask{The one idea, and it covers all three lectures.} Every model here is
built from the same two moves: compute a set of \emph{weights} that say how much
each thing should contribute, and take a weighted average. Attention does it over
positions in a sequence. A VAE does it over a distribution of latent codes.
Diffusion does it over time steps of a noising process. The differences are in
what is being averaged and what determines the weights.
\end{keybox}

\section{Attention}

\subsection{Tokens are a matrix}

A sequence of $T$ tokens, each embedded as a $d$-dimensional vector, is a matrix
$\boldsymbol{X}$ of shape $(T, d)$. Everything in a transformer operates on that
matrix. Note what has been given up compared to an RNN: there is no order in a
matrix's rows unless you put it there, which is why positional encodings exist.

\subsection{Q, K, V are three views of the same input}

\begin{lstlisting}
Q = X @ Wq        # (T, d_k)   what each position is looking for
K = X @ Wk        # (T, d_k)   what each position offers
V = X @ Wv        # (T, d_v)   what each position would contribute
\end{lstlisting}

All three are linear projections of the \emph{same} $\boldsymbol{X}$ --- that is
what makes it \emph{self}-attention. The names are a retrieval metaphor: a query
is matched against keys, and the matching keys' values are returned.

\subsection{Scaled dot-product attention, in four lines}

\[
\mathrm{Attention}(\boldsymbol{Q},\boldsymbol{K},\boldsymbol{V})
= \mathrm{softmax}\!\left(\frac{\boldsymbol{Q}\boldsymbol{K}^{\!\top}}{\sqrt{d_k}}\right)\boldsymbol{V}
\]

\begin{lstlisting}
scores = Q @ K.transpose(-2, -1) / math.sqrt(d_k)    # (T, T)
if mask is not None:
    scores = scores.masked_fill(mask == 0, float('-inf'))
A = torch.softmax(scores, dim=-1)                    # (T, T), rows sum to 1
out = A @ V                                          # (T, d_v)
\end{lstlisting}

Track the shapes, because they are the whole thing:

\begin{center}
\begin{tabular}{@{}lll@{}}
\toprule
Step & Shape & Meaning \\
\midrule
$\boldsymbol{Q}\boldsymbol{K}^{\!\top}$ & $(T, T)$ & every position against every position \\
after softmax & $(T, T)$ & each \emph{row} is a distribution over positions \\
$\boldsymbol{A}\boldsymbol{V}$ & $(T, d_v)$ & each position: a weighted average of all values \\
\bottomrule
\end{tabular}
\end{center}

\begin{keybox}
\ask{Three details that are not decoration.}
\textbf{The $\sqrt{d_k}$}: dot products of $d_k$-dimensional vectors grow like
$\sqrt{d_k}$, and a softmax of large numbers saturates to a one-hot vector with
almost no gradient. Remove it and training fails for large $d_k$.
\textbf{\texttt{dim=-1}}: the softmax runs across \emph{keys}, so each row is a
distribution over what to attend to. Softmax over the wrong axis produces
nonsense that runs perfectly.
\textbf{The mask is $-\infty$, not $0$}: it must be applied \emph{before} the
softmax, because zeroing afterwards leaves rows that no longer sum to one.
\end{keybox}

\subsection{Multi-head, which is only a reshape}

Split the $d$ dimensions into $H$ groups, run the same operation on each, and
concatenate the results. It lets different heads attend to different things.

\begin{lstlisting}
# (N, T, d) -> (N, H, T, d//H)
Q = Q.view(N, T, H, d // H).transpose(1, 2)
# ... attention exactly as above, batched over the H axis ...
out = out.transpose(1, 2).reshape(N, T, d)     # concatenate the heads back
\end{lstlisting}

The \texttt{transpose(1, 2)} before and after is the only subtle part. It puts
the head axis next to the batch axis so that the matrix multiplies treat heads
as independent problems.

\subsection{Reading an attention map, carefully}

The $(T, T)$ matrix $\boldsymbol{A}$ is tempting to plot and interpret. Three
cautions, all of which have caused published overclaims:

\begin{itemize}
  \item There are $H$ heads per layer and many layers, so ``the attention map''
        is a choice among dozens, and they disagree.
  \item Attention weights are \emph{not} feature importances. A position can be
        attended to heavily and contribute little, because what it contributes
        is its \emph{value} vector, not itself.
  \item The residual stream carries information around the attention block
        entirely, so a position with zero attention is not disconnected.
\end{itemize}

\begin{keybox}
\ask{The one-sentence version.} An attention map tells you where information was
\emph{routed}, not what the model \emph{used}. Recitation 9's whole discipline
applies here too.
\end{keybox}

\section{Autoregressive generation}

Training a language model is next-token prediction with teacher forcing ---
exactly the mechanism from Recitation 10, at scale. Generation is a loop:

\begin{lstlisting}
for _ in range(n_new):
    logits = model(idx)[:, -1, :] / temperature      # last position only
    if top_k is not None:
        v, _ = torch.topk(logits, top_k)
        logits[logits < v[:, [-1]]] = float('-inf')
    probs = torch.softmax(logits, dim=-1)
    nxt   = torch.multinomial(probs, num_samples=1)  # SAMPLE, do not argmax
    idx   = torch.cat([idx, nxt], dim=1)
\end{lstlisting}

\begin{center}
\begin{tabular}{@{}ll@{}}
\toprule
Knob & What it does \\
\midrule
\texttt{temperature} $\to 0$ & approaches greedy: repetitive, often degenerate \\
\texttt{temperature} $= 1$ & samples from the model's own distribution \\
\texttt{temperature} $> 1$ & flatter, more surprising, less coherent \\
\texttt{top\_k} & restricts to the $k$ most likely tokens before sampling \\
\bottomrule
\end{tabular}
\end{center}

\begin{keybox}
\ask{Sampling is not a detail of the model.} Two very different outputs from the
same weights differ only in these two numbers, so any claim about ``what the
model does'' has to state them. A surprising fraction of published prompting
results are partly results about a temperature setting.
\end{keybox}

A useful measurement, and one you can do on a laptop: \textbf{surprisal}, the
negative log probability the model assigned to the token that actually occurred.
It is the quantity that correlates with human reading times, which is the whole
basis of using language models as models of language processing.

\begin{lstlisting}
logp = torch.log_softmax(model(idx).logits, dim=-1)
surprisal = -logp[0, :-1].gather(1, idx[0, 1:, None]) / math.log(2)   # in bits
\end{lstlisting}

\section{Generative models, two objectives}

\subsection{VAE: the reparameterization trick and the ELBO}

The encoder outputs a \emph{distribution} over codes rather than a code. You
cannot backpropagate through ``draw a sample'', so you rewrite the sample as a
deterministic function of the parameters plus fixed noise:

\[
\boldsymbol{z} = \boldsymbol{\mu} + \boldsymbol{\sigma}\odot\boldsymbol{\varepsilon},
\qquad \boldsymbol{\varepsilon}\sim\mathcal{N}(0, \boldsymbol{I}).
\]

Now the randomness sits in $\boldsymbol{\varepsilon}$, which needs no gradient,
and $\boldsymbol{\mu}$ and $\boldsymbol{\sigma}$ are on the differentiable path.

\begin{lstlisting}
mu, logvar = encoder(x).chunk(2, dim=1)      # predict log-variance, not variance
std = torch.exp(0.5 * logvar)
z   = mu + std * torch.randn_like(std)       # the reparameterization trick
xhat = decoder(z)

recon = F.mse_loss(xhat, x, reduction='sum')
kl    = -0.5 * torch.sum(1 + logvar - mu.pow(2) - logvar.exp())
loss  = recon + beta * kl
\end{lstlisting}

The network predicts \texttt{logvar} rather than the variance so that its output
is unconstrained --- a variance must be positive, and $\exp$ enforces that for
free. The two loss terms pull against each other: reconstruction wants codes
that are distinct and informative, the KL wants them close to a standard normal
so that the space is smooth and samplable. \texttt{beta} sets the exchange rate,
and the whole character of the model follows from it.

\begin{keybox}
\ask{The failure mode with a name.} If the KL term dominates --- large
\texttt{beta}, or a decoder powerful enough to reconstruct without help --- the
encoder collapses to the prior, $\boldsymbol{\mu}\to0$ and
$\boldsymbol{\sigma}\to1$, and the latent code carries \emph{nothing}. Samples
still look fine because the decoder learned the data on its own. This is
posterior collapse, and the diagnostic is to check whether the KL term goes to
zero.
\end{keybox}

\subsection{Diffusion: one closed form and a five-line loss}

Adding Gaussian noise repeatedly has a closed form, so you never simulate the
forward process step by step:

\[
\boldsymbol{x}_t = \sqrt{\bar\alpha_t}\,\boldsymbol{x}_0 +
\sqrt{1-\bar\alpha_t}\,\boldsymbol{\varepsilon}.
\]

That is the whole reason training is cheap: pick a random $t$, jump straight
there, and ask the network to name the noise.

\begin{lstlisting}
t   = torch.randint(0, T, (x0.shape[0],))
eps = torch.randn_like(x0)
xt  = abar[t].sqrt()[:, None] * x0 + (1 - abar[t]).sqrt()[:, None] * eps
loss = F.mse_loss(model(xt, t), eps)          # predict the noise you added
\end{lstlisting}

The training objective is a mean squared error. That is all it is. The
complexity of diffusion models lives entirely in the \emph{sampling} loop, which
runs the learned denoiser backwards many times, and in the architecture of the
denoiser --- not in the loss.

\section{Exercises}

\ask{1.} You delete the $\sqrt{d_k}$ from scaled dot-product attention. Training
works at $d_k = 8$ and fails at $d_k = 512$. Why?

\ask{2.} You apply a causal mask by multiplying the attention weights by 0
\emph{after} the softmax. What is wrong?

\ask{3.} You take softmax over \texttt{dim=-2} instead of \texttt{dim=-1}. Does
anything error? What does the model now compute?

\ask{4.} A colleague shows an attention map and says it proves the model bases
its answer on the word ``not''. Give two objections.

\ask{5.} Generation at \texttt{temperature=0.1} produces the same sentence
repeated forever. Explain in terms of the sampling loop.

\ask{6.} Your VAE's KL term drops to almost zero within a few epochs, and
reconstructions look good. What has happened, and is the model useful?

\ask{7.} Diffusion training is a plain MSE. Where did all the complexity go?

\subsection*{Answers}

\ask{1.} Dot products of $d_k$-dimensional vectors have magnitude growing like
$\sqrt{d_k}$. At $d_k = 512$ the unscaled scores are large, the softmax saturates
to nearly one-hot, and its gradient is nearly zero, so nothing trains. At
$d_k = 8$ the scores are small enough that it hardly matters.

\ask{2.} The softmax has already distributed weight to the masked positions, so
zeroing afterwards leaves rows that do not sum to one --- and the unmasked
positions get less than their share. Mask with $-\infty$ \emph{before} the
softmax.

\ask{3.} Nothing errors: the matrix is $(T, T)$, so both axes are valid. You now
normalize down the columns, so each row is no longer a distribution over
positions and the output is a weighted combination with weights that do not sum
to one. It trains, badly, with no diagnostic.

\ask{4.} (i) There are many heads and many layers; ``the attention map'' is one
choice among dozens that disagree. (ii) Attention weight is not importance --- a
position contributes its \emph{value} vector, which may be near zero, and the
residual stream carries information around the block regardless.

\ask{5.} At low temperature the softmax is nearly one-hot, so sampling is
effectively \texttt{argmax}. The model deterministically picks its most likely
continuation, enters a state it has seen before, and cycles. Raising the
temperature or using \texttt{top\_k} with sampling breaks the loop.

\ask{6.} Posterior collapse: the encoder has matched the prior and the latent
code carries no information, with the decoder reconstructing on its own. It is
not useful as a \emph{representation} --- $\boldsymbol{z}$ tells you nothing about
$\boldsymbol{x}$ --- even though samples may look fine. Lower \texttt{beta}, or
weaken the decoder.

\ask{7.} Into the sampling loop, which runs the learned denoiser backwards
across many steps, and into the denoiser architecture. The training objective is
genuinely just ``predict the noise you added'', which is why diffusion models are
much easier to train than GANs.

\section*{Cheat sheet}
\addcontentsline{toc}{section}{Cheat sheet}

\begin{center}
\begin{tabular}{@{}p{0.44\textwidth}p{0.5\textwidth}@{}}
\toprule
\textbf{You want} & \textbf{You write} \\
\midrule
attention scores          & \texttt{Q @ K.transpose(-2,-1) / sqrt(d\_k)} \\
a causal mask             & \texttt{masked\_fill(mask == 0, float('-inf'))} \emph{before} softmax \\
the softmax axis          & \texttt{dim=-1} --- over keys, always \\
split into heads          & \texttt{view(N,T,H,d//H).transpose(1,2)} \\
sample a token            & \texttt{torch.multinomial(probs, 1)} --- not \texttt{argmax} \\
surprisal in bits         & \texttt{-log\_softmax(...).gather(...) / log(2)} \\
the reparameterization    & \texttt{mu + torch.exp(0.5*logvar) * randn\_like(mu)} \\
the KL term               & \texttt{-0.5 * sum(1 + logvar - mu**2 - logvar.exp())} \\
a diffusion training step & \texttt{mse\_loss(model(xt, t), eps)} \\
\bottomrule
\end{tabular}
\end{center}

\end{document}
